% FILE: artifacts/mixture-input-encoding_poster_ccn2026/parts/fonts.tex
% ==================================================================================================
% Typography (LuaLaTeX + fontspec). Two families realise the plan's hierarchy:
%   - Raleway          : display geometric sans -> title (Black) and headings (SemiBold/Bold)
%   - Source Sans 3  : neutral sans          -> body, captions, labels, tables
% Fonts are resolved by name from the TeX Live tree (raleway, sourcesans packages).

\usepackage{fontspec}
\defaultfontfeatures{Ligatures=TeX}

% --- Body / default family: Source Sans 3 (sans) ---
% The ROMAN font is Latin Modern (serif) so that math alphabets (\mathrm, \mathbf) and the
% math "operators" font are serif, consistent with the serif math symbols. The document default
% is still sans (familydefault = sfdefault), so body text is Source Sans 3.
\setmainfont{Latin Modern Roman}
\setsansfont{SourceSans3-Regular.otf}[BoldFont=SourceSans3-Bold.otf,
  ItalicFont=SourceSans3-RegularIt.otf, BoldItalicFont=SourceSans3-BoldIt.otf]
\renewcommand{\familydefault}{\sfdefault}

% --- Display family: Raleway (weights addressed explicitly) ---
% Title uses a lighter Medium weight (Black reads too heavy at poster scale).
\newfontfamily\titlefont{Raleway Medium}
\newfontfamily\headingfont{Raleway SemiBold}[
  BoldFont       = {Raleway Bold},
  ItalicFont     = {Raleway SemiBold Italic},
  BoldItalicFont = {Raleway Bold Italic}
]

% --- Supporting body weights ---
\newfontfamily\bodysemibold{SourceSans3-Semibold.otf}

% --- Mono family for section-heading kickers / metadata (per tests/claude-design-layout.tex) ---
\newfontfamily\kickerfont{Source Code Pro}

% --- Math text font: the body sans, so descriptive WORDS inside a formula (amplification, projection,
% comp. plane, ...) match the body prose. Symbols, vectors, matrices and the "mix" indices keep the
% math font. Used via \mlab (macros.tex). ---
\newfontfamily\mathtextfont{SourceSans3-Regular.otf}

% --- MATH SIZES, ONE DECLARATION PER STEP OF THE TYPE SCALE ---------------------------------------
% \DeclareMathSizes{<text>}{<math>}{<script>}{<scriptscript>} fires when the TEXT size is selected --
% inside a TikZ node too. EVERY step of the scale in macros.tex that can carry math needs a line here;
% a step without one falls back to LaTeX's defaults and drifts out of the system SILENTLY.
%
% MATH = 1.125 x TEXT. The body is Source Sans 3, the math Latin Modern (serif): its smaller x-height
% makes a symbol read small beside the sans at the same nominal size, and the 12.5% bump levels them.
%
% SCRIPT / SCRIPTSCRIPT keep the conventional 0.7 / 0.5 -- deliberate, since the floors in macros.tex
% govern TEXT, not indices. KNOWN CONSEQUENCE: at body size the "mix" index lands at 23 pt, just under
% the 24 pt figure floor and the smallest mark on the poster. If it fails at viewing distance, raise
% the last two numbers here; nothing else changes.
\DeclareMathSizes{24}{27}{19}{14}      % \figticksize  -- figure floor
\DeclareMathSizes{30}{34}{24}{17}      % \figannotsize
\DeclareMathSizes{36}{40}{28}{20}      % \notesize / \figcapsize -- text floor
\DeclareMathSizes{40}{45}{32}{23}      % \bodysize / \bodycompact -- incl. the two move (a) formulae
\DeclareMathSizes{44}{50}{35}{25}      % \cardtitlesize
\DeclareMathSizes{52}{58}{41}{29}      % \headingsize
% The brace readings under the move (a) formulae do NOT follow these sizes -- they are sans, pinned by
% \brlab (macros.tex), because a gloss is words, not algebra.
